\section{Controlled Relaxation Protocol}
\label{sec:relaxation_protocol}

The structural root of the Wamura problem is that evidence
\emph{against} a risk claim requires exercise of the restricted
capability, which the restriction prevents.  Controlled relaxation
breaks this deadlock by creating bounded, monitored exercise
windows that generate genuine Bayesian evidence while bounding
maximum damage.

\subsection{The test object}

\begin{definition}[Controlled Relaxation Test]
\label{def:cr_test}
A \emph{controlled relaxation test}
$T = (R, \Xirestrict, \scope, \testdur, \damagebound)$ for a
risk-restriction pair $(R, \Xirestrict)$ consists of:
\begin{itemize}
\item The risk-restriction pair $(R, \Xirestrict)$ being tested;
\item A \emph{test scope} $\scope \subseteq P$: a bounded sub-poset
  within which $\Xirestrict$ is temporarily lifted.  The scope
  determines maximum $\volL$ exposure
  (Section~\ref{sec:damage_guarantees});
\item A \emph{test duration} $\testdur \in \mathbb{N}$: the number
  of time steps for which $\Xirestrict$ is lifted within $\scope$;
  after $\testdur$ steps, $\Xirestrict$ is reinstated regardless of
  outcome;
\item A \emph{damage tolerance} $\damagebound > 0$: the maximum
  acceptable $\volL$-contraction from the test if the risk
  materializes.  $\damagebound$ is set by the meta-policy
  (Section~\ref{sec:meta_policy}) or by the user as a safety bound.
\end{itemize}
\end{definition}

\subsection{Protocol phases}

\paragraph{Phase 1: Scope selection.}
Identify a sub-poset $\scope$ such that
\begin{itemize}
\item[(i)] the predicted pathway $\Pathway$ is exercisable within
  $\scope$ (the prerequisite capabilities are available);
\item[(ii)] the maximum $\volL$-contraction from any exercise
  pathway within $\scope$ is bounded by $\damagebound$
  (Theorem~\ref{thm:damage_bound});
\item[(iii)] no other controlled relaxation test is currently
  running on a scope $\scope'$ that overlaps with $\scope$
  (isolation, to prevent cascading test failures).
\end{itemize}

\paragraph{Phase 2: Relaxation.}
Within scope $\scope$, lift $\Xirestrict$ for $\testdur$ time steps.
During the test window:
\begin{itemize}
\item Monitor for the predicted contraction pathway using the
  multi-channel detection architecture
  of~\citep[Table~1]{lasser2026scm}: capability gating, causal
  attribution, behavioural residuals, SPRT detection,
  correlational co-occurrence.
\item Record all observations in the test log with timestamps and
  channel attributions (formalized in
  Section~\ref{sec:detection_sensitivity}).
\item If the detection channel reports $\volL$-contraction within
  $\scope$ exceeding $\damagebound$ at any step, immediately
  reinstate $\Xirestrict$ and terminate the test early (safety
  cutoff).  The cutoff and the Bayesian-update observable $O_n$
  (Phase~3) read from the same noisy detector; a cutoff event is
  therefore a detection event in the sense of $O_n = 1$, not a
  separate exact-contraction channel.
\end{itemize}

\paragraph{Phase 3: Observation and Bayesian update.}
We separate the latent event (pathway activation) from the
observed event (detection alarm).  The model is a \emph{simple
binary hypothesis test} with fixed likelihoods throughout:

\begin{itemize}
\item $A_n \in \{0, 1\}$: indicator that the pathway activates
  during test $n$ (latent; not directly observed);
\item $O_n \in \{0, 1\}$: indicator that contraction is detected
  during test $n$ (observed);
\item $H_0$: the pathway is harmless---activation probability
  $\theta = 0$ (pathway never activates even when exercised);
\item $H_1$: the pathway is dangerous---activation probability
  $\theta = \thetaone$ for a fixed $\thetaone \in (0, 1]$
  (pathway activates with known probability $\thetaone$ when
  exercised).
\end{itemize}

The choice of $\thetaone$ parametrizes the alternative hypothesis
as a fixed point in $(0, 1]$.  The conservative default is
$\thetaone = 1$ (``if the pathway is dangerous at all, it
activates every time the capability is exercised''), which
produces the strongest evidence per test.  A more nuanced
$\thetaone < 1$ reduces evidence per test but may be more
realistic for stochastically activating pathways.  Under either
choice, $\thetaone$ is fixed across the test sequence and is
\emph{not} a free parameter in the likelihood.

\paragraph{Observation model.}
The detector is characterised by sensitivity and specificity:
\begin{align*}
\Pr(O_n = 1 \mid A_n = 1) &= 1 - \beta(\testdur)
  \qquad (\text{true-positive rate}),\\
\Pr(O_n = 1 \mid A_n = 0) &= \alphafp
  \qquad (\text{false-positive rate}),
\end{align*}
where $\beta(\testdur)$ is the false-negative rate over the test
window (characterised in Section~\ref{sec:detection_sensitivity})
and $\alphafp \in [0, 1)$ is the per-test false-positive rate of
the aggregate detector.

\paragraph{Fixed per-hypothesis likelihoods.}
Marginalizing over $A_n$ under each hypothesis gives, under $H_1$:
\begin{align}
\label{eq:lik_h1}
\Pr(O_n = 1 \mid H_1)
  &= \thetaone (1 - \beta) + (1 - \thetaone) \alphafp, \\
\Pr(O_n = 0 \mid H_1)
  &= \thetaone \beta + (1 - \thetaone) (1 - \alphafp),
\nonumber
\end{align}
and under $H_0$:
\begin{align}
\label{eq:lik_h0}
\Pr(O_n = 1 \mid H_0) &= \alphafp,\\
\Pr(O_n = 0 \mid H_0) &= 1 - \alphafp.
\nonumber
\end{align}
These are \emph{fixed numbers} determined by $(\thetaone, \beta,
\alphafp)$, not functions of a free parameter.  This is what makes
the binary test well-defined.

\paragraph{Likelihood ratios.}
Define the per-test likelihood ratios
\begin{align}
\label{eq:L_defs}
\Lone &= \frac{\Pr(O_n = 1 \mid H_1)}{\Pr(O_n = 1 \mid H_0)}
   = \frac{\thetaone (1 - \beta) + (1 - \thetaone) \alphafp}{\alphafp},\\
\Lzero &= \frac{\Pr(O_n = 0 \mid H_1)}{\Pr(O_n = 0 \mid H_0)}
   = \frac{\thetaone \beta + (1 - \thetaone) (1 - \alphafp)}{1 - \alphafp}.
\nonumber
\end{align}
Under the identifiability assumption
(Assumption~\ref{ass:T1} in Section~\ref{sec:convergence}) we will
have $\Lzero < 1$ and $\Lone > 1$; both are fixed scalars
depending only on $(\thetaone, \beta, \alphafp)$.

\paragraph{Posterior updates.}
For outcome $O_n = 1$ (contraction detected),
\begin{equation}
\label{eq:update_pos}
\pR(t^+) = \frac{\pR(t) \, \Lone}{\pR(t) \, \Lone + (1 - \pR(t))},
\qquad
\logit \pR(t^+) = \logit \pR(t) + \log \Lone.
\end{equation}
For outcome $O_n = 0$ (no contraction detected),
\begin{equation}
\label{eq:update_neg}
\pR(t^+) = \frac{\pR(t) \, \Lzero}{\pR(t) \, \Lzero + (1 - \pR(t))},
\qquad
\logit \pR(t^+) = \logit \pR(t) + \log \Lzero.
\end{equation}
Because $\Lzero < 1$ under Assumption~\ref{ass:T1}, each
no-contraction observation pushes the posterior toward $H_0$.
This is what breaks the Wamura asymmetry.

\begin{remark}[Contrast with the restricted regime]
\label{rem:regime_contrast}
Under the restricted regime
(Proposition~\ref{prop:asym_evidence}), $A_n = 0$ with certainty
(the restriction prevents activation), so both outcome likelihoods
collapse to the common value and the effective $\Lzero = \Lone =
1$: the posterior is unchanged.  Under controlled relaxation, the
capability is exercised (Assumption~\ref{ass:T5} in
Section~\ref{sec:convergence}), activation can occur, and the
detector can observe it---making both outcomes informative and
breaking the ratchet.
\end{remark}

\begin{remark}[Continuous extension]
\label{rem:continuous_extension}
The binary formulation ($H_0$ versus $H_1$ with fixed $\thetaone$)
is sufficient for the operational question ``should this
restriction be maintained?'' and avoids the dimensional mismatch
between a binary posterior and a continuous activation rate.  A
continuous extension placing a prior $\pi(\theta)$ on $[0,1]$
could estimate the exact activation rate but requires latent-state
augmentation for the posterior update: the latent $A_n$ makes the
model non-conjugate with a Beta prior, unlike the standard
Beta--Binomial.  We defer this extension to future work; the
binary formulation is what the Wamura problem requires.
\end{remark}

\subsection{Connection to the phase boundary}
\label{sec:connection_phase_boundary}

\citet{lasser2026phase} show that feedback loop~3 (over-calibrated
restriction $\to$ discovery avoidance $\to$ structural ignorance)
has a degenerate phase boundary when the self-correction rate
$\rS = 0$ for that
loop~\citep[Proposition~4]{lasser2026phase}.  The self-correction
rate is zero precisely because the Wamura asymmetry prevents
evidence accumulation on the restricted dimension.  Controlled
relaxation makes $\rS > 0$ by injecting exercise-regime
observations into the evidence stream: the log-odds update
of~\eqref{eq:update_pos}--\eqref{eq:update_neg} has nonzero
expected drift toward the true hypothesis under
Assumption~\ref{ass:T1}, and that drift is the restoring force
required by the design criterion
of~\citep[Theorem~2]{lasser2026phase} to lift Loop~3 off the
phase-boundary degeneracy.
